\section{Threat Model}
\label{sec:threat}

\subsection{System Roles and Deployment Contracts}

The defender deploys an online-updated controller $\pi_{\theta_t}$ and a
certificate mechanism; a severe-attack-aware estimator or filter returns a
plausible-state set $X_{\mathcal A}(h_t)$ and certified action kernel
$K_\Gamma(h_t)$.  A trusted state anchor, integrity-protected fallback
controller, or external fail-safe may also be available, but each is an explicit architectural
assumption and cost.
Unless a formal safety result is stated, \emph{trusted} and
\emph{integrity-protected} mean outside the attacker's write authority; they do
not mean that a controller or handoff is verified safe.

The evaluation instantiates four contracts.  \emph{Permanent raw release}
removes runtime authority after a one-time five-step raw-policy test: protected
commissioning is followed by snapshot export and promotion to a lower-latency
control target, which avoids permanent computation, intervention, and
dependence on a backup service.  \emph{Commit admission} certifies a candidate
snapshot on a declared state set and horizon before raw release.
The \emph{raw-policy lookahead switcher} repeats the same check online and can
transfer control to an integrity-protected LQR fallback.  The check supplies
evidence about the raw policy, not about LQR recovery from the handoff state;
the latter is a separate handoff obligation.  The \emph{one-step action
filter} applies a current-state action filter and serves as a negative
control.  A sound filter that remains online instead certifies the composition
$F\circ\pi_{\theta^+}$.  Each name denotes a distinct evidence scope and
authority contract, and Section~\ref{sec:study} measures their tradeoffs;
Appendix~\ref{sec:appendix-delta} gives a contract taxonomy.

\subsection{Split-Trust Learning Pipeline}

The split-trust model separates a high-integrity runtime-assurance plane
from an asynchronous learning-data plane.  The runtime plane owns the
authenticated plant state, action selector, actuator interface,
integrity-protected fallback, and certificate outcome.  The learning plane consumes recorded
transitions and scalar rewards through telemetry, a historian, a replay
buffer, or an external key-performance-indicator (KPI) interface.  OT guidance
recognizes this separation and the exchange of historian data across zones
\cite{nist_ot_security_2023,mitre_data_historian}.  The benchmark models the
resulting interface without assuming a particular product.

\begin{table}[t]
\centering
\caption{Evaluated write boundary and the read authority the evaluated attack
consumes.  ``Integrity-protected'' means outside attacker write authority, not
formally verified or certified safe.}
\label{tab:trust-boundary}
\footnotesize
\begin{tabular}{@{}L{0.23\columnwidth}L{0.26\columnwidth}L{0.19\columnwidth}L{0.13\columnwidth}@{}}
\toprule
Object & Defender assumption & Write access & Read access\\
\midrule
Runtime state, action, plant & Authenticated control path & No write & Logged copy\\
Fallback, selector, certificate & Integrity-protected runtime code/state & No write & No read\\
Transitions and learner code & Integrity-protected after collection & No write & Read\\
Reward record at ingestion & No origin/semantic check & Bounded scalar write & Read\\
Parameters and optimizer state & Written only by updater & No write & Read\\
Held-out states and release rule & Fixed, honest evaluation & No write or selection & No read\\
\bottomrule
\end{tabular}
\end{table}

Table~\ref{tab:trust-boundary} states the evaluated boundary in both
directions.  The split-trust argument constrains write authority; the read
column records what the evaluated attack rule consumes, including the learner
state against which its edits are computed.  The attacker is therefore
write-narrow but knowledge-rich, rather than weak in every capability
(Section~\ref{sec:limits}).  The attack begins
after compromise of a principal with bounded reward-field write access, which permits post-computation, pre-update mutation by a replay writer
or broker.  A compromised reward producer is equivalent when it has the same
bounded output authority; in that case source authentication establishes
origin but not reward semantics.

\emph{Applicability condition.}  The learning-data plane is deliberately
reachable: historians, ingestion brokers, replay buffers, and KPI services
exist so that offline consumers can read and enrich operational records, and
OT guidance treats cross-zone historian exchange as an expected data
flow~\cite{nist_ot_security_2023,mitre_data_historian}.  A principal on that
path holds bounded write authority over a scalar field long before it holds
any authority over the control path, which is why we grant the attacker
nothing else.  The channel is only interesting when the runtime plane cannot
recompute the reward from authenticated transitions.  That is the case when
reward depends on delayed quality labels, operator feedback, manual
disposition codes, or an external KPI the runtime plane does not observe.  It
is not the case for rewards that are deterministic functions of authenticated
state and action: there, trusted recomputation removes the channel outright,
as our frozen integrity audit confirms (Section~\ref{sec:study}).  Both
simulation platforms expose deterministic rewards, so trusted recomputation is
available in the experiments.  We use them to isolate the consequences of the
reward-ingestion interface, not as evidence that these benchmarks inherently
lack recomputation.  The deployment attack surface is limited to otherwise
equivalent interfaces whose reward semantics cannot be reconstructed inside
the high-integrity plane.

\subsection{Attack Channels and Capabilities}

The general history model permits two attack channels.  Observation FDI
$a_t$ changes the state semantics of the history and enlarges
$X_{\mathcal A}(h_t)$.  Update-data corruption $b_t$ modifies learning signals
and can redirect $U(\theta_t,h_t)$ with trusted state observations.  Either
channel can alter the next action map.

The end-to-end learner experiments use the narrower reward-record channel.
The white-box, batch-aware attacker observes the current learner and
adaptation batch, but not the held-out deployment-state choice.  It modifies
only scalar rewards before an ordinary policy-gradient update, subject to a
declared per-step $L_\infty$ budget.  It cannot write any object marked ``No write'' in
Table~\ref{tab:trust-boundary}.  The adaptation
selector observes each proposed action before the plant step and retains
physical authority during data collection.

\subsection{Dependency-Aware Handoff Evidence}
\label{sec:evidence-transfer}

The contracts above differ in their evidence obligations.  We record an
evidence item as
\begin{equation}
    E=(\varphi,\mathcal D; s,a,p,i,\tau).
\end{equation}
Here $\varphi$ is the claim supported by the evidence and $\mathcal D$ is the
set of premises on which that support depends.  The remaining fields describe
scope.  The \emph{subject} $s$ is the raw policy $\pi_\theta$ or the assured
composition $F\circ\pi_\theta$.  The \emph{authority} $a$ identifies the
resident module that constrains applied actions, or its absence.  The
\emph{provenance} $p$ records the process that supplied learning inputs.  The
\emph{information} $i$ describes the state, observation, belief, or filtered
action available at decision time.  The \emph{horizon} $\tau$ is the interval
over which the evidence bounds behavior.

The record is an audit device, not a proof calculus.  A changed field prompts
reassessment but does not automatically invalidate evidence.  Evidence $E$
is insufficient for a post-handoff contract when that contract no longer
entails some premise in $\mathcal D$.  Conversely, a field outside
$\mathcal D$ may change without affecting the claim.  For example, a direct
verification of the final raw snapshot can remain valid even if its derivation
does not assume trusted training provenance.  By contrast, a training trace of
the shielded composition does not establish raw-policy behavior, and a finite
raw-policy test does not establish behavior beyond its declared horizon.

The controlled study contains two handoffs with different missing premises.
Protected adaptation followed by raw execution changes the acting composition.
The five-step raw-policy test supplies new evidence for the released policy and
the absence of runtime intervention, but only over five steps.  Extending that
evidence to 120-step execution changes the required horizon.  The clean cohort
therefore measures horizon extrapolation after release.  The raw-policy
lookahead switcher creates a later handoff from the monitored policy to the LQR
fallback.  Its trigger predicts raw-policy behavior but does not establish LQR
recovery from the transfer state.

Prior systems illustrate both sides of this audit.  Independent transfer
evidence can discharge changed premises before deployment
\cite{hsu_sim_lab_real_2023}; Simplex mechanisms retain a decision module and
use controller-specific switching conditions
\cite{seto_simplex_1998,phan_neural_simplex_2020,damare_bb_simplex_2025}.
Carr et al. supply a published transition from shielded bootstrap to raw-policy
authority, but not the five-step release contract used in our controlled
study~\cite{carr_partial_shielding_2023}.  Our case study adds adversarial
learning-input provenance to their retirement lifecycle.

\subsection{Integrity Assumptions and Escape Conditions}

The vulnerable configuration supplies neither origin integrity for the reward
record nor trusted reward recomputation.  Binding a reward, transition,
action, producer identity, and time to an authenticated append-only record
blocks post-ingestion mutation, and recomputation by a trusted principal
completely blocks the evaluated channel whenever reward is a deterministic
function of authenticated transitions.  Under the applicability condition
above, the system instead needs a trusted producer, redundant validation, or a
detector; a signature applied to an already malicious value is insufficient.

These controls act at different boundaries: provenance and recomputation
remove the upstream precondition, a task-aware detector can reject suspicious
batches under its distributional and semantic assumptions, commit admission
restricts which learned snapshot can be released, and the raw-policy lookahead
switcher adds a downstream handoff whose safety depends on the LQR fallback and
the state at which it receives control.  In the frozen
audit, trusted recomputation achieved true- and false-positive rates (TPR/FPR) of
1.000/0 and simple task-aware checks achieved TPR above 0.90 and FPR at most
0.036 against the attack rule we first evaluated.  Only recomputation closes
the evaluated channel: an attacker that respects the task-aware predicates by
construction retains deployment harm, and the fallback handoff can still be
followed by a violation (Section~\ref{sec:constrained}).  Those checks bound one
attack rule rather than the channel or the downstream handoff.

\subsection{Security and Availability Goals}

We separate four attacker or failure objectives.  An \emph{unsafe certified
action} occurs when the mechanism accepts an applied input that is unsafe for
some state in $X_{\mathcal A}(h_t)$.  A \emph{stale certified policy} occurs
when the current applied input is filtered but the accepted raw updated action
map leaves the kernel.  \emph{Uncertified learning} occurs when a nontrivial
update continues after certificate rejection.  \emph{Forced freeze} occurs
when the attacker expands the plausible set or redirects the update until
sound learning must stop or re-anchor.  The defender also measures the
availability metrics of Section~\ref{sec:lifecycle-attacks}; useful learning
requires paired task improvement over always-freeze in addition to physical
safety.

Severe-attack filters, secure estimators, and fault-tolerant control-barrier
function (CBF) methods can preserve safety when their assumptions keep
$K_\Gamma(h_t)$ nonempty
\cite{fawzi_secure_2014,tan_safety_2024,zhang_safe_2022}.  Out of scope are
service exploitation, compromise of trusted recomputation, and attacks on the
authenticated runtime state or fallback controller; Section~\ref{sec:limits}
states the remaining validity boundaries.
